% !TeX root = ../MQF_Manuale_Master.tex
% Capitoli/MQF_Cap_16_Python_Programmazione_Stocastica.tex


\chapter{Applicazione in Python: programmazione stocastica}

\label{cap:python-programmazione-stocastica}

I Capitoli 14 e 15 hanno costruito il quadro teorico necessario per
rappresentare decisioni finanziarie assunte in presenza di incertezza. Il
Capitolo 14 ha introdotto la programmazione stocastica a due stadi, distinguendo
le decisioni assunte prima dell'osservazione dello scenario dalle decisioni di
ricorso che possono essere adattate all'informazione successivamente
disponibile. Il Capitolo 15 ha sviluppato il confronto tra soluzione
stocastica, soluzione basata sui valori attesi e soluzione con informazione
perfetta, introducendo il valore della soluzione stocastica e il valore
atteso dell'informazione perfetta.

Questo capitolo compie il passaggio successivo. L'obiettivo non \`e introdurre
una nuova classe di modelli, ma rendere operativa in Python la struttura
stocastica gi\`a formulata, traducendo decisioni di primo stadio, variabili di
ricorso, scenari, probabilit\`a e vincoli nella forma richiesta dal solver.

Rispetto ai problemi deterministici affrontati nei Capitoli 11--13, la
differenza essenziale \`e che una stessa decisione iniziale deve essere
compatibile con pi\`u possibili realizzazioni future. Le variabili di ricorso
possono cambiare da uno scenario all'altro, mentre la decisione assunta prima
che l'incertezza si risolva deve rimanere unica. La struttura computazionale
deve quindi preservare esplicitamente la distinzione

\[
\text{decisione ex ante}
\longrightarrow
\text{scenario}
\longrightarrow
\text{decisione di ricorso}.
\]


Il filo conduttore del capitolo \`e un Caso Aula ispirato alla crisi di
Silicon Valley Bank del 2023. Si considera un problema stilizzato di
Asset--Liability Management nel quale una banca deve scegliere l'allocazione
iniziale tra attivit\`a caratterizzate da rendimenti e gradi di liquidabilit\`a
differenti. Dopo la decisione iniziale, uno tra pi\`u scenari di liquidit\`a
pu\`o realizzarsi e rendere necessario il ricorso a funding di emergenza.

Nel passaggio all'implementazione, non verr\`{a} riprodotta integralmente
la struttura multiperiodale del caso SVB--ALM utilizzata nei Capitoli 14
e 15. Per isolare gli aspetti computazionali propri della formulazione
a due stadi, il Caso Aula adotter\`{a} una versione ridotta del modello:
l'incertezza sar\`{a} rappresentata mediante un insieme finito di scenari
di liquidit\`{a} e, per ciascuno scenario, il ricorso sar\`{a} descritto
da un unico blocco di decisioni. Rimangono invariati i principi essenziali
del modello a due stadi: una decisione iniziale comune agli scenari,
ricorso scenario-specifico e valutazione mediante probabilit\`{a}.

Il problema permette di osservare direttamente il ruolo della
non anticipativit\`a. L'allocazione iniziale deve essere scelta prima di
conoscere lo scenario futuro e rimane quindi comune a tutte le realizzazioni.
Le variabili di funding e le eventuali giacenze residue possono invece
adattarsi allo scenario osservato. L'incertezza non viene quindi trattata
risolvendo separatamente pi\`u problemi deterministici, ma costruendo un unico
programma lineare nel quale decisione iniziale e ricorso scenario-specifico
sono collegati.

Il capitolo introduce inoltre un piccolo focus Python dedicato alla
traduzione di un programma lineare nella forma utilizzata da
\texttt{scipy.optimize.linprog}. La funzione era gi\`a stata utilizzata nel
Capitolo 13 per il problema ALM deterministico; qui viene ripresa per mostrare
come l'ordinamento delle variabili, la funzione obiettivo, i vincoli di
uguaglianza, le disuguaglianze e i bound debbano essere organizzati quando il
vettore decisionale comprende contemporaneamente variabili comuni e variabili
specifiche dei diversi scenari.

Una volta risolto il problema stocastico, l'analisi non si esaurisce nella
determinazione della soluzione ottima. La stessa struttura computazionale
viene utilizzata per costruire e confrontare tre differenti regimi
decisionali. La soluzione stocastica considera esplicitamente tutti gli
scenari e le loro probabilit\`a; il problema Expected Value sostituisce i
parametri incerti con i rispettivi valori medi; il problema Wait-and-See
suppone invece che lo scenario sia noto prima della decisione iniziale.

Il confronto richiede particolare attenzione. La soluzione ottenuta dal
problema medio non pu\`o essere confrontata direttamente con la soluzione
stocastica utilizzando soltanto il valore della funzione obiettivo del
problema deterministico. La decisione Expected Value deve essere fissata e
rivalutata negli scenari originari. Solo questa rivalutazione produce il
benchmark necessario per calcolare il valore della soluzione stocastica.

Il percorso quantitativo del capitolo pu\`o essere sintetizzato come

\[
\begin{array}{c}
\text{modello stocastico}
\\[3pt]
\downarrow
\\[3pt]
\text{forma computazionale}
\\[3pt]
\downarrow
\\[3pt]
\text{SP, EV, WS}
\\[3pt]
\downarrow
\\[3pt]
\text{VSS ed EVPI}
\\[3pt]
\downarrow
\\[3pt]
\text{controllo e interpretazione}.
\end{array}
\]

Particolare attenzione sar\`a dedicata ai controlli. Una soluzione dichiarata
ottima dal solver non garantisce, da sola, che il problema sia stato tradotto
correttamente. Devono essere verificati l'ordinamento delle variabili, i
vincoli di bilancio, i bound, la fattibilit\`a delle decisioni fissate nei
diversi scenari e la coerenza dei benchmark utilizzati nel confronto.

Questi controlli assumono rilievo anche nell'uso dell'intelligenza
artificiale. Una risposta formalmente plausibile pu\`o contenere valori non
coerenti con gli output del notebook, modificare implicitamente la
specificazione del problema o trascurare un vincolo di fattibilit\`a. Il
metodo di lavoro rimane quello introdotto nel Capitolo 4 per i capitoli
applicativi e viene qui assunto come acquisito: la Scheda Caso definisce la
specifica, il flusso logico-teorico precede la costruzione del notebook e ogni
risultato deve essere sottoposto a controllo prima di essere interpretato.

Accanto al Caso Aula, viene proposto un Caso TakeHome ispirato alla crisi dei
fondi Liability-Driven Investment nel Regno Unito nel 2022. In questo caso,
la decisione iniziale riguarda il bilanciamento tra buffer di liquidit\`a ed
esposizione LDI, mentre il ricorso comprende liquidit\`a esterna e
deleveraging in presenza di margin call. Cambiano il contesto finanziario e
la struttura del ricorso, ma rimane invariata la domanda fondamentale:
scegliere una decisione ex ante che resti compatibile con gli scenari futuri
e valutarne correttamente il costo in presenza di informazione incompleta.

L'obiettivo complessivo del capitolo \`e quindi rendere operativo il passaggio
dalla programmazione stocastica teorica alla decisione finanziaria
computabile e verificabile. Python traduce il modello in un problema risolto
dal solver; i benchmark distinguono differenti regimi informativi; i
controlli stabiliscono se i risultati ottenuti possono essere effettivamente
interpretati sul piano economico.



\section{Python per risolvere programmi lineari con \texttt{linprog}}

Nel Capitolo 13, \texttt{scipy.optimize.linprog} \`e stato utilizzato per
tradurre e risolvere un problema lineare deterministico. Nel presente
capitolo, lo stesso solver viene applicato a un problema di programmazione
stocastica a due stadi.

La forma generale rimane quella gi\`a nota:

\[
\min_{v}\; c^{\prime}v
\]

\noindent soggetto a

\[
A_{\mathrm{ub}}v\leq b_{\mathrm{ub}},
\qquad
A_{\mathrm{eq}}v=b_{\mathrm{eq}},
\]

\noindent insieme ai bound sulle singole variabili.

La novit\`a non riguarda quindi il solver, ma la struttura del vettore
decisionale e dei vincoli.

\subsection{Variabili comuni e variabili scenario-specifiche}

In un problema a due stadi, il vettore decisionale comprende una parte comune
a tutti gli scenari e una parte che pu\`o adattarsi allo scenario osservato:

\[
v=
\left(
x,
y_1,
y_2,
\ldots,
y_S
\right)^{\prime}.
\]

Il vettore \(x\) contiene le decisioni assunte prima che lo scenario sia noto;
il vettore \(y_s\) contiene invece le variabili di ricorso associate allo
scenario \(s\).

Nel Caso Aula, le decisioni iniziali sono

\[
x_1,x_2,x_3,x_4,
\]

\noindent mentre per ciascuno scenario
\(s\in\{N,P,R\}\) compaiono una variabile di funding \(u_s\) e una variabile di giacenza
residua \(g_s\).

Un possibile ordinamento \`e

\[
v=
\left(
x_1,x_2,x_3,x_4,
u_N,u_P,u_R,
g_N,g_P,g_R
\right)^{\prime}.
\]

L'ordine scelto deve rimanere invariato nella funzione obiettivo, nelle
matrici dei vincoli e nei bound.

\subsection{La struttura a blocchi dei vincoli}

Il vincolo sulle risorse iniziali riguarda soltanto le decisioni di primo
stadio:

\[
x_1+x_2+x_3+x_4=A_0.
\]

I vincoli di liquidit\`a sono invece specifici dei diversi scenari. Per lo
scenario \(s\),

\[
\sum_{j=1}^{4}\alpha_{js}x_j
+
u_s
=
d_s+g_s.
\]

Nella matrice \(A_{\mathrm{eq}}\), le colonne associate a \(x\) compaiono
quindi nei vincoli di tutti gli scenari, mentre le colonne relative a
\(u_s\) e \(g_s\) compaiono soltanto nella riga dello scenario
corrispondente.

La struttura computazionale riflette direttamente la sequenza

\[
x
\longrightarrow
s
\longrightarrow
(u_s,g_s).
\]

Il fatto che \(x\) compaia una sola volta nel vettore decisionale rende
esplicita la non anticipativit\`a della decisione iniziale.

\subsection{Funzione obiettivo e chiamata al solver}

Se il problema economico \`e formulato come massimizzazione,

\[
\max_v\; c_{\Pi}^{\prime}v,
\]

\noindent \texttt{linprog} viene utilizzato con

\[
c=-c_{\Pi}.
\]

La chiamata essenziale rimane

\begin{verbatim}
result = linprog(
    c,
    A_ub=A_ub,
    b_ub=b_ub,
    A_eq=A_eq,
    b_eq=b_eq,
    bounds=bound,
    method="highs"
)
\end{verbatim}

Il vettore ottimo \`e contenuto in

\[
\texttt{result.x},
\]

\noindent mentre, nel caso di inversione del segno della funzione obiettivo,

\[
z^*
=
-\texttt{result.fun}.
\]

\subsection{Il controllo essenziale}

Nel problema stocastico, il controllo pi\`u importante riguarda la
corrispondenza tra posizione delle variabili e significato economico.

Devono essere verificati almeno:

\begin{itemize}
\item l'ordinamento del vettore decisionale;

\item la corretta associazione tra righe dei vincoli e scenari;

\item i bound assegnati alle singole variabili;

\item la presenza di un'unica decisione \(x\) comune a tutti gli scenari;

\item la corretta estrazione delle variabili di ricorso dalla soluzione.
\end{itemize}

Il fatto che

\[
\texttt{result.success=True}
\]

\noindent indica soltanto che il solver ha risolto il problema numerico che
gli \`e stato fornito. Non garantisce che la formulazione implementata
coincida con il modello economico assegnato.

Nel seguito, lo stesso strumento verr\`a utilizzato per costruire e
confrontare la soluzione stocastica, il problema Expected Value e i problemi
Wait-and-See.



\section{Dalla programmazione stocastica al problema computazionale}

I Capitoli 14 e 15 hanno formulato il problema stocastico a due stadi
distinguendo una decisione iniziale \(x\), assunta prima dell'osservazione
dello scenario, e una decisione di ricorso \(y_s\), che pu\`o dipendere dallo
scenario \(s\).

La struttura essenziale \`e

\[
x
\longrightarrow
s
\longrightarrow
y_s.
\]

Nel problema computazionale, questa sequenza deve essere trasformata in un
unico sistema di variabili e vincoli nel quale la decisione \(x\) rimane
comune a tutti gli scenari, mentre ciascun \(y_s\) compare soltanto nel blocco
corrispondente.

Se gli scenari appartengono all'insieme

\[
\mathcal{S}
=
\{1,\ldots,S\},
\]

\noindent con probabilit\`a \(p_s\), una formulazione generale pu\`o essere
scritta come

% \[
% \max_{x,\{y_s\}}
% \left\{
% c^{\prime}x
% +
% \sum_{s\in\mathcal{S}}
% p_s q_s^{\prime}y_s
% \right\},
% \]

\[
\begin{aligned}
\max_{x,\{y_s\}_{s\in\mathcal S}}\quad&
c^{\prime}x+
\sum_{s\in\mathcal S}p_s q_s^{\prime}y_s\\
\text{soggetto a}\quad&
x\in\mathcal X_0,\\
&\text{vincoli di ricorso nello scenario }s,
\qquad s\in\mathcal S.
\end{aligned}
\]

%\noindent soggetta ai vincoli di primo stadio e ai vincoli di ricorso specifici dei diversi scenari.

La rappresentazione computazionale non richiede quindi di risolvere un
problema separato per ciascuno scenario. Tutti i blocchi vengono raccolti
nello stesso programma lineare e condividono la medesima decisione iniziale.

\subsection{La soluzione stocastica}

La soluzione stocastica, indicata con \(x^{SP}\), viene ottenuta utilizzando
congiuntamente tutti gli scenari e le rispettive probabilit\`a.

Il corrispondente valore ottimo \`e

\[
z^{SP}.
\]

Dal punto di vista computazionale, questo \`e il problema nel quale la matrice
dei vincoli contiene simultaneamente tutti i blocchi scenario-specifici e la
funzione obiettivo pondera i costi o i benefici di ricorso mediante \(p_s\).

La soluzione \(x^{SP}\) rappresenta quindi la decisione ex ante scelta tenendo
conto dell'intera distribuzione degli scenari.

\subsection{Il problema Expected Value e la rivalutazione}

Nel problema Expected Value, i parametri incerti vengono sostituiti dai
rispettivi valori medi.
Indicando genericamente con \(\xi_s\) il vettore dei parametri dello scenario,

\[
\bar{\xi}
=
\sum_{s\in\mathcal{S}}
p_s \xi_s.
\]

Si ottiene cos\`i un problema deterministico, la cui soluzione viene indicata
con

\[
x^{EV}.
\]

Il valore ottimo di questo problema medio non costituisce per\`o il termine di
confronto corretto con \(z^{SP}\).

La decisione \(x^{EV}\) deve essere fissata e rivalutata nel problema
stocastico originario. In ciascuno scenario vengono quindi ottimizzate
soltanto le variabili di ricorso, mantenendo invariata la decisione di primo
stadio.

Il valore atteso cos\`i ottenuto \`e indicato con

\[
z^{EV}.
\]

Per un problema di massimizzazione, il valore della soluzione stocastica \`e

\[
VSS
=
z^{SP}
-
z^{EV}.
\]

Il confronto misura il beneficio derivante dall'avere considerato
esplicitamente l'incertezza nella scelta della decisione iniziale.

\subsection{Il problema Wait-and-See}

Il benchmark Wait-and-See modifica invece l'informazione disponibile al
momento della decisione.

Per ciascuno scenario \(s\), si risolve un problema deterministico assumendo
che lo scenario sia gi\`a noto prima della scelta di \(x\). La decisione
iniziale pu\`o quindi cambiare da uno scenario all'altro.

Indicando con \(z_s^{WS}\) il valore ottimo nello scenario \(s\),

\[
z^{WS}
=
\sum_{s\in\mathcal{S}}
p_s z_s^{WS}.
\]

Per un problema di massimizzazione,

\[
EVPI
=
z^{WS}
-
z^{SP}.
\]

Il valore \(EVPI\) misura quindi il vantaggio ottenibile disponendo di
informazione perfetta prima della decisione iniziale.

\subsection{Tre problemi, una stessa struttura}

I benchmark utilizzano la stessa base modellistica, ma differiscono per il
modo in cui viene trattata l'informazione:

\[
\begin{array}{ccl}
SP
&:&
\text{decisione unica prima dello scenario},
\\[3pt]
EV
&:&
\text{decisione sul problema costruito con parametri medi},
\\[3pt]

WS
&:&
\text{decisione diversa per ciascuno scenario noto}.
\end{array}
\]

Quando tutti i benchmark risultano ben definiti e fattibili, per un problema
di massimizzazione deve valere

\[
z^{EV}
\leq
z^{SP}
\leq
z^{WS}.
\]

Questa relazione costituisce anche un controllo della procedura
computazionale.

Nel Caso Aula, i quattro problemi verranno costruiti a partire dalla stessa
struttura di Asset--Liability Management. Ci\`o che cambier\`a non sar\`a il
solver utilizzato, ma l'informazione disponibile al momento della decisione e
il modo in cui la decisione iniziale viene mantenuta comune, fissata, o resa
scenario-specifica.


\section{Caso Aula: SVB--ALM 2023 --- problema, modello e flusso risolutivo}

Il Caso Aula considera un problema stilizzato di Asset--Liability Management
ispirato alle tensioni di liquidit\`a emerse durante la crisi di Silicon
Valley Bank nel 2023.

L'obiettivo non \`e ricostruire storicamente il bilancio della banca, ma
rappresentare in forma semplificata un problema decisionale nel quale una
allocazione iniziale deve essere scelta prima di conoscere il fabbisogno di
liquidit\`a futuro.

La banca dispone inizialmente di risorse pari a

\[
A_0=100
\]

\noindent e deve distribuirle tra quattro categorie di attivit\`a.

Indichiamo con

\[
x=
\left(
x_1,x_2,x_3,x_4
\right)^{\prime}
\]

\noindent il vettore delle allocazioni iniziali.

Le quattro componenti rappresentano, rispettivamente:

\begin{itemize}
\item \(x_1\): liquidit\`a e riserve;

\item \(x_2\): titoli a breve termine;

\item \(x_3\): titoli a lunga scadenza;

\item \(x_4\): prestiti e altre attivit\`a meno liquide.
\end{itemize}

I rendimenti assegnati sono

\[
c=
\left(
0.010,\,
0.025,\,
0.040,\,
0.050
\right)^{\prime}.
\]

L'allocazione iniziale deve soddisfare

\[
x_1+x_2+x_3+x_4=100,
\]

\[
x_j\geq 0,
\]

\noindent insieme ai limiti

\[
x_3\leq 70,
\qquad
x_4\leq 70.
\]

In assenza di incertezza, il problema tenderebbe a privilegiare le attivit\`a
con rendimento pi\`u elevato. La presenza di possibili fabbisogni di
liquidit\`a modifica tuttavia il confronto, perch\'e le diverse categorie di
attivo non sono ugualmente liquidabili nei diversi scenari.

\subsection{Scenari di liquidit\`a e capacit\`a di realizzo}

Dopo la scelta di \(x\), si realizza uno tra tre scenari:

\[
\mathcal{S}
=
\{N,P,R\},
\]

\noindent dove \(N\) rappresenta una condizione normale, \(P\) una condizione
di pressione e \(R\) uno scenario di run.

Le probabilit\`a assegnate sono

\[
p_N=0.60,
\qquad
p_P=0.25,
\qquad
p_R=0.15.
\]

Il fabbisogno di liquidit\`a associato ai tre scenari \`e

\[
d_N=15,
\qquad
d_P=35,
\qquad
d_R=65.
\]

La quantit\`a di liquidit\`a ottenibile da ciascuna attivit\`a dipende dallo
scenario attraverso i coefficienti \(\alpha_{js}\).

La matrice dei coefficienti di liquidabilit\`a \`e

\[
\begin{array}{c|ccc}
 & N & P & R\\
\hline
x_1 & 1.00 & 1.00 & 1.00\\
x_2 & 0.98 & 0.95 & 0.90\\
x_3 & 0.85 & 0.70 & 0.50\\
x_4 & 0.60 & 0.40 & 0.20
\end{array}
\]

La struttura rappresenta una progressiva riduzione della capacit\`a di
trasformare le attivit\`a in liquidit\`a quando le condizioni di mercato
peggiorano.

La liquidit\`a ottenibile dall'allocazione iniziale nello scenario \(s\)
\`e quindi

\[
\sum_{j=1}^{4}\alpha_{js}x_j.
\]

Lo stesso portafoglio pu\`o pertanto generare quantit\`a di liquidit\`a molto
diverse nei tre scenari.

\subsection{Le decisioni di ricorso}

Dopo l'osservazione dello scenario, la banca pu\`o ricorrere a funding di
emergenza.

Indichiamo con

\[
u_s\geq 0
\]

\noindent il funding utilizzato nello scenario \(s\) e con

\[
g_s\geq 0
\]

\noindent l'eventuale liquidit\`a residua dopo avere soddisfatto il
fabbisogno.

Per ciascuno scenario deve valere il bilancio

\[
\sum_{j=1}^{4}\alpha_{js}x_j
+
u_s
=
d_s+g_s.
\]

La stessa equazione svolge due funzioni.

Se la liquidit\`a generata dall'attivo \`e inferiore al fabbisogno, il
funding \(u_s\) colma la differenza.

Se invece la liquidit\`a disponibile supera il fabbisogno, l'eccedenza viene
registrata in \(g_s\).

Il funding ha un costo scenario-specifico. I coefficienti assegnati sono

\[
\kappa_N=0.01,
\qquad
\kappa_P=0.06,
\qquad
\kappa_R=0.25.
\]

Il costo associato al ricorso nello scenario \(s\) \`e quindi

\[
\kappa_s u_s.
\]

Il costo cresce fortemente nello scenario pi\`u avverso. La composizione
dell'attivo deve pertanto essere valutata non soltanto in base al rendimento,
ma anche in base alla capacit\`a di contenere il ricorso a funding costoso.

\subsection{La funzione obiettivo}

Il rendimento iniziale del portafoglio \`e

\[
\sum_{j=1}^{4}c_jx_j.
\]

Il costo atteso del funding \`e invece

\[
\sum_{s\in\mathcal{S}}
p_s\kappa_su_s.
\]

Il problema stocastico consiste quindi nel massimizzare

\[
\sum_{j=1}^{4}c_jx_j
-
\sum_{s\in\mathcal{S}}
p_s\kappa_su_s.
\]

La formulazione completa \`e

\[
\max_{x,u,g}
\left\{
\sum_{j=1}^{4}c_jx_j
-
\sum_{s\in\mathcal{S}}
p_s\kappa_su_s
\right\}
\]

\noindent soggetta a

\[
x_1+x_2+x_3+x_4=100,
\]

\[
\sum_{j=1}^{4}\alpha_{js}x_j
+
u_s
=
d_s+g_s,
\qquad
s\in\mathcal{S},
\]

\[
x_j\geq 0,
\qquad
u_s\geq 0,
\qquad
g_s\geq 0,
\]

\[
x_3\leq 70,
\qquad
x_4\leq 70.
\]

Il funding di emergenza non presenta un limite superiore. Per ogni
allocazione ammissibile \(x\), un eventuale deficit di liquidit\`a pu\`o
quindi essere compensato aumentando \(u_s\). Nel caso considerato, il
problema dispone pertanto di ricorso fattibile nei tre scenari.

\subsection{La struttura decisionale del caso}

La sequenza del problema pu\`o essere rappresentata come

\[
\begin{array}{c}
\text{scelta dell'allocazione }x
\\[3pt]
\downarrow
\\[3pt]
\text{realizzazione dello scenario }s
\\[3pt]
\downarrow
\\[3pt]
\text{liquidit\`a ottenibile dall'attivo}
\\[3pt]
\downarrow
\\[3pt]
\text{funding }u_s\text{ e giacenza }g_s
\\[3pt]
\downarrow
\\[3pt]
\text{costo di ricorso}
\\[3pt]
\downarrow
\\[3pt]
\text{valore atteso della politica}.
\end{array}
\]

La decisione \(x\) precede l'osservazione dello scenario e deve quindi
rimanere unica. Le variabili \(u_s\) e \(g_s\) vengono invece determinate
dopo l'osservazione di \(s\) e possono cambiare da uno scenario all'altro.

Il problema non consiste dunque nel trovare l'allocazione migliore in
ciascuno scenario separatamente. Consiste nel trovare una sola allocazione
iniziale che tenga conto congiuntamente di tutti gli scenari possibili.

\subsection{Dal problema al flusso risolutivo}

La Scheda Caso assegna variabili, parametri, vincoli e benchmark da
calcolare. Per organizzare la soluzione computazionale, questi elementi
devono essere disposti secondo le loro dipendenze logiche.

Il flusso risolutivo pu\`o essere sintetizzato in sei passaggi.

\begin{center}
\small
\begin{tabular}{p{0.8cm}p{3.0cm}p{4.1cm}p{3.6cm}}
\textbf{Passo} &
\textbf{Finalit\`a} &
\textbf{Operazione} &
\textbf{Output o controllo}
\\ \hline

1 &
Definire la decisione ex ante &
Organizzare rendimenti, risorse iniziali e limiti di allocazione &
Spazio ammissibile delle decisioni \(x\)
\\

2 &
Rappresentare l'incertezza &
Organizzare probabilit\`a, fabbisogni e coefficienti di liquidabilit\`a &
Coerenza degli scenari e somma delle probabilit\`a
\\

3 &
Costruire il ricorso &
Associare a ogni scenario \(u_s\), \(g_s\) e il relativo bilancio &
Fattibilit\`a dei bilanci scenario-specifici
\\

4 &
Risolvere il problema stocastico &
Ottimizzare congiuntamente \(x\) e le variabili di ricorso &
\(x^{SP}\), ricorsi ottimi e \(z^{SP}\)
\\

5 &
Costruire i benchmark informativi &
Risolvere EV, rivalutare \(x^{EV}\) e risolvere i problemi WS &
\(z^{EV}\) e \(z^{WS}\)
\\

6 &
Confrontare le politiche &
Calcolare VSS ed EVPI e verificare l'ordinamento dei benchmark &
Interpretazione economica e controlli finali
\end{tabular}
\end{center}

Il flusso mantiene separati tre livelli.

Il primo riguarda la decisione iniziale e la sua esposizione ai diversi
scenari. Il secondo riguarda la risposta disponibile dopo l'osservazione
dello scenario. Il terzo riguarda il confronto tra differenti strutture
informative.

La sezione successiva tradurr\`a questa struttura nel notebook, costruendo in
sequenza la soluzione stocastica, il benchmark Expected Value con la relativa
rivalutazione e il benchmark Wait-and-See.


\section{Dal flusso al notebook: SP, EV e WS}

Una volta definita la struttura del problema, il notebook pu\`o essere
organizzato attorno ai tre benchmark introdotti nel Capitolo 15: soluzione
stocastica, Expected Value con rivalutazione e Wait-and-See.

I tre problemi condividono la stessa base economica, ma differiscono per
l'informazione disponibile al momento della decisione iniziale.

\subsection{La soluzione stocastica}

Nel problema stocastico, la decisione iniziale

\[
x=
\left(
x_1,x_2,x_3,x_4
\right)^{\prime}
\]

\noindent viene scelta prima dell'osservazione dello scenario e rimane quindi
comune a \(N\), \(P\) e \(R\).

Le variabili di ricorso

\[
u_N,u_P,u_R,
\qquad
g_N,g_P,g_R
\]

\noindent possono invece adattarsi alla condizione di liquidit\`a
effettivamente realizzata.

Con l'ordinamento

\[
v=
\left(
x_1,x_2,x_3,x_4,
u_N,u_P,u_R,
g_N,g_P,g_R
\right)^{\prime},
\]

\noindent il notebook costruisce il vettore della funzione obiettivo, i
vincoli di uguaglianza e i bound coerentemente con la formulazione della
sezione precedente.

La soluzione ottenuta \`e

\[
x^{SP}
=
(0,30,70,0)^{\prime}.
\]

L'intero ammontare iniziale viene quindi allocato tra titoli a breve e titoli
a lunga scadenza, con il limite massimo su \(x_3\) completamente utilizzato.

Il rendimento lordo dell'allocazione \`e

\[
0.025\cdot 30
+
0.040\cdot 70
=
3.55.
\]

Nei tre scenari, il ricorso a funding di emergenza risulta

\[
u_N^{SP}=0,
\qquad
u_P^{SP}=0,
\qquad
u_R^{SP}=3.
\]

Il funding viene quindi utilizzato soltanto nello scenario di run.

Il costo di funding nello scenario \(R\) \`e

\[
\kappa_R u_R^{SP}
=
0.25\cdot 3
=
0.75.
\]

Questa quantit\`a non coincide con il contributo dello scenario al costo
atteso. Poich\'e

\[
p_R=0.15,
\]

\noindent il contributo ponderato \`e

\[
p_R\kappa_Ru_R^{SP}
=
0.15\cdot 0.75
=
0.1125.
\]

Il costo atteso complessivo del funding \`e quindi

\[
0.1125,
\]

\noindent e il valore della soluzione stocastica risulta

\[
z^{SP}
=
3.55-0.1125
=
3.4375.
\]

Il risultato mostra gi\`a il compromesso fondamentale del problema. La
soluzione non massimizza semplicemente il rendimento iniziale delle
attivit\`a, ma sceglie una composizione che limita il costo di ricorso negli
scenari sfavorevoli.

\subsection{Expected Value e rivalutazione della decisione}

Il problema Expected Value sostituisce i parametri incerti con i rispettivi
valori medi.

Il fabbisogno medio \`e

\[
\bar d
=
\sum_{s\in\mathcal{S}}p_s d_s
=
27.5.
\]

I coefficienti medi di liquidabilit\`a sono

\[
\bar\alpha
=
\left(
1,\,
0.9605,\,
0.76,\,
0.49
\right)^{\prime},
\]

\noindent mentre il costo medio del funding \`e

\[
\bar\kappa
=
0.0585.
\]

Il problema deterministico costruito con questi parametri produce

\[
x^{EV}
=
(0,0,30,70)^{\prime}.
\]

Rispetto alla soluzione stocastica, l'allocazione si sposta verso le
attivit\`a con rendimento pi\`u elevato.

Il valore ottimo del problema medio \`e 4.7.

Questo valore, tuttavia, non \`e \(z^{EV}\) e non deve essere confrontato
direttamente con \(z^{SP}\).

Il problema medio utilizza infatti una sola configurazione costruita con
parametri medi. Per valutare correttamente la decisione \(x^{EV}\), occorre
fissarla e riportarla negli scenari originari.

La rivalutazione mantiene quindi

\[
x=x^{EV}
\]

\noindent e ottimizza soltanto le variabili di ricorso nei tre scenari.

Nel caso normale e in quello di pressione non \`e necessario ricorrere a
funding:

\[
u_N^{EV}=0,
\qquad
u_P^{EV}=0.
\]

Nello scenario di run, invece,

\[
u_R^{EV}=36.
\]

La maggiore esposizione alle attivit\`a meno liquide rende quindi la
decisione basata sui valori medi sensibilmente pi\`u costosa nello scenario
estremo.

Il valore atteso della decisione Expected Value rivalutata negli scenari
originari \`e

\[
z^{EV}
=
3.35.
\]

Il valore della soluzione stocastica \`e pertanto

\[
VSS
=
z^{SP}-z^{EV}
=
3.4375-3.35
=
0.0875.
\]

Il valore numerico del \(VSS\) \`e contenuto, ma non deve essere interpretato
come prova dell'equivalenza economica tra le due politiche.

Le allocazioni

\[
x^{SP}
=
(0,30,70,0)^{\prime}
\]

\noindent e

\[
x^{EV}
=
(0,0,30,70)^{\prime}
\]

\noindent sono infatti profondamente diverse, cos\`i come il ricorso al funding
nello scenario pi\`u avverso.

Il confronto mostra quindi che una differenza modesta nel valore atteso della
funzione obiettivo pu\`o accompagnarsi a una differenza rilevante nella
struttura di liquidit\`a e nell'esposizione allo stress.

\subsection{Il benchmark Wait-and-See}

Il problema Wait-and-See assume che lo scenario sia noto prima della scelta
dell'allocazione iniziale.

La decisione \(x\) non deve quindi essere comune a tutti gli scenari, ma pu\`o
essere scelta separatamente per \(N\), \(P\) e \(R\).

Negli scenari normale e di pressione, la soluzione \`e

\[
x_N^{WS}
=
x_P^{WS}
=
(0,0,30,70)^{\prime}.
\]

In queste condizioni, l'informazione anticipata consente di privilegiare le
attivit\`a con rendimento pi\`u elevato senza sostenere un costo di ricorso
sufficiente a modificare l'allocazione.

Nello scenario di run, la soluzione cambia invece in modo marcato:

\[
x_R^{WS}
=
(0,37.5,62.5,0)^{\prime}.
\]

La conoscenza anticipata dello scenario estremo induce quindi una maggiore
allocazione verso titoli a breve e l'abbandono delle attivit\`a meno liquide.

I valori ottimi scenario-specifici sono

\[
z_N^{WS}=4.7,
\qquad
z_P^{WS}=4.7,
\qquad
z_R^{WS}=3.4375.
\]

Il valore atteso Wait-and-See \`e pertanto

\[
z^{WS}
=
0.60\cdot 4.7
+
0.25\cdot 4.7
+
0.15\cdot 3.4375
=
4.510625.
\]

Il valore atteso dell'informazione perfetta risulta

\[
EVPI
=
z^{WS}-z^{SP}
=
4.510625-3.4375
=
1.073125.
\]

Nel caso considerato, il valore dell'informazione perfetta \`e quindi molto
pi\`u elevato del valore della soluzione stocastica rispetto alla politica
basata sui valori medi.

I due indicatori misurano per\`o fenomeni differenti.

Il \(VSS\) confronta la decisione stocastica con una decisione costruita
sostituendo l'incertezza con i valori medi.

L'\(EVPI\) confronta invece la decisione stocastica con una situazione
irrealistica nella quale lo scenario futuro \`e noto prima dell'allocazione.

\subsection{Confrontare i benchmark}

I risultati possono essere riassunti come

\[
\begin{array}{c|c|c}
\text{Benchmark}
&
\text{Allocazione iniziale}
&
\text{Valore}
\\ \hline
EV
&
(0,0,30,70)
&
3.35
\\
SP
&
(0,30,70,0)
&
3.4375
\\
WS
&
\text{scenario-specifica}
&
4.510625
\end{array}
\]

e verificano l'ordinamento

\[
z^{EV}
\leq
z^{SP}
\leq
z^{WS}.
\]

Nel caso specifico,

\[
3.35
\leq
3.4375
\leq
4.510625.
\]

La relazione costituisce un controllo sintetico, ma non esaurisce la verifica
del notebook. La coerenza dei benchmark richiede anche che:

\begin{itemize}
\item \(x^{SP}\) sia unica e comune ai tre scenari;

\item \(x^{EV}\) venga mantenuta fissa durante la rivalutazione;

\item le decisioni \(x_s^{WS}\) possano variare tra gli scenari;

\item i costi di ricorso siano distinti dai relativi contributi ponderati alla
funzione obiettivo;

\item ogni soluzione rispetti i vincoli e i bound del problema corrispondente.
\end{itemize}

Il confronto finale non riguarda quindi soltanto tre valori ottimi. Riguarda
tre differenti strutture informative e il modo in cui ciascuna di esse
modifica la decisione iniziale e il ricorso successivo.


\section{Controlli, fattibilit\`a e interpretazione economica}

La soluzione di un problema di programmazione stocastica non termina quando il
solver restituisce un vettore ottimo. Prima dell'interpretazione economica,
occorre verificare che il problema implementato coincida con quello definito
dalla Scheda Caso e che i diversi benchmark siano stati costruiti in modo
coerente.

Nel Caso Aula, i controlli riguardano tre livelli distinti: la traduzione del
modello, la fattibilit\`a delle soluzioni e la coerenza dei confronti tra SP,
EV e WS.

\subsection{Controllare la traduzione del modello}

Il primo livello riguarda la corrispondenza tra formulazione economica e
rappresentazione matriciale.

Devono essere verificati:

\begin{itemize}
\item l'ordinamento delle variabili nel vettore decisionale;

\item la corretta associazione dei coefficienti della funzione obiettivo;

\item la presenza di un'unica decisione iniziale \(x\) nel problema
stocastico;

\item la corretta collocazione delle variabili \(u_s\) e \(g_s\) nei vincoli
dello scenario corrispondente;

\item la coerenza dei bound con i limiti assegnati nella Scheda Caso.
\end{itemize}

Un errore in uno di questi elementi pu\`o produrre un problema perfettamente
risolvibile dal solver ma economicamente diverso da quello assegnato.

Per questo motivo,

\[
\texttt{result.success=True}
\]

\noindent rappresenta soltanto un controllo sul problema numerico fornito al
solver, non sulla correttezza della sua costruzione.

\subsection{Controllare la fattibilit\`a}

Una volta ottenuta la soluzione, i vincoli devono essere ricostruiti
utilizzando i valori ottimi.

Per ciascuno scenario \(s\), il bilancio

\[
\sum_{j=1}^{4}\alpha_{js}x_j
+
u_s
=
d_s+g_s
\]

\noindent deve essere soddisfatto, insieme ai vincoli

\[
x_j\geq 0,
\qquad
u_s\geq 0,
\qquad
g_s\geq 0,
\]

\[
x_3\leq 70,
\qquad
x_4\leq 70.
\]

La stessa verifica deve essere effettuata per i benchmark costruiti
successivamente.

Nel problema EV, in particolare, la decisione \(x^{EV}\) deve rimanere
fissata. La rivalutazione riguarda soltanto le variabili di ricorso. Se la
decisione iniziale venisse nuovamente ottimizzata, il problema non
rappresenterebbe pi\`u la valutazione della politica Expected Value.

Nel benchmark WS, avviene invece l'opposto: la decisione iniziale pu\`o essere
ottimizzata separatamente in ciascuno scenario, poich\'e l'informazione viene
supposta disponibile prima della scelta.

La distinzione tra questi casi non riguarda quindi una semplice opzione
computazionale, ma la struttura informativa del problema.

\subsection{Controllare i valori economici}

Un ulteriore controllo riguarda la costruzione della funzione obiettivo.

Nel Caso Aula, il costo di funding nello scenario \(s\) \`e

\[
\kappa_su_s,
\]

\noindent mentre il suo contributo al costo atteso \`e

\[
p_s\kappa_su_s.
\]

Le due quantit\`a non devono essere confuse.

Analogamente, il valore ottimo del problema Expected Value non coincide con
\(z^{EV}\). Il primo \`e ottenuto nel problema deterministico costruito con
parametri medi; il secondo misura il valore della decisione \(x^{EV}\) quando
viene riportata negli scenari originari.

Quando i tre benchmark sono correttamente definiti e fattibili, deve inoltre
valere

\[
z^{EV}
\leq
z^{SP}
\leq
z^{WS}.
\]

Nel Caso Aula,

\[
3.35
\leq
3.4375
\leq
4.510625.
\]

La relazione fornisce un controllo sintetico della coerenza complessiva del
confronto, ma deve essere utilizzata dopo avere verificato separatamente le
soluzioni da cui i tre valori derivano.

\subsection{Interpretare la soluzione}

La soluzione stocastica produce

\[
x^{SP}
=
(0,30,70,0)^{\prime},
\]

\noindent mentre il problema Expected Value conduce a

\[
x^{EV}
=
(0,0,30,70)^{\prime}.
\]

La seconda allocazione offre un rendimento lordo pi\`u elevato, ma concentra
una quota maggiore nelle attivit\`a meno liquide. Quando la stessa decisione
viene rivalutata nello scenario di run, il ricorso al funding aumenta
fortemente.

La soluzione stocastica rinuncia quindi a parte del rendimento potenziale per
ridurre l'esposizione al costo di ricorso negli scenari avversi.

Il valore

\[
VSS=0.0875
\]

\noindent pu\`o apparire modesto se osservato isolatamente. Tuttavia, il
confronto tra le allocazioni mostra che una piccola differenza nella funzione
obiettivo attesa pu\`o accompagnarsi a una differenza rilevante nella
struttura finanziaria della politica adottata.

Il \(VSS\) misura infatti una differenza di valore atteso, non una distanza
tra portafogli n\'e una misura diretta della vulnerabilit\`a di liquidit\`a.

Il valore

\[
EVPI=1.073125
\]

\noindent risulta invece molto pi\`u elevato. L'informazione perfetta
permetterebbe di scegliere un'allocazione diversa prima di ciascuno scenario e
di evitare il compromesso imposto dalla decisione ex ante.

Questo benchmark non rappresenta una strategia realisticamente disponibile,
ma quantifica il costo economico dell'incertezza residua al momento della
decisione.


\subsection{Controllo numerico e responsabilit\`a interpretativa}

Quando il notebook viene sviluppato con il supporto dell'intelligenza
artificiale, i risultati proposti devono essere verificati rispetto alla
Scheda Caso, al modello e agli output effettivamente prodotti.

Il criterio operativo rimane

\[
\text{output}
\longrightarrow
\text{controllo}
\longrightarrow
\text{interpretazione}.
\]

L'IA pu\`o assistere la verifica e suggerire possibili letture economiche, ma
non sostituisce il controllo della coerenza tra specificazione, codice e
risultati.

Nel Caso Aula, il punto centrale \`e che l'incertezza modifica la composizione
ottima dell'attivo: considerare esplicitamente gli scenari conduce a una
politica meno esposta al fabbisogno di funding nelle condizioni pi\`u
avverse.


\section{Caso TakeHome: UK LDI 2022 --- buffer, margin call e fattibilit\`a del ricorso}

Il Caso TakeHome trasferisce la struttura della programmazione stocastica a
due stadi a un problema ispirato alla crisi dei fondi Liability-Driven
Investment nel Regno Unito nel 2022.

La decisione iniziale riguarda la ripartizione di risorse pari a

\[
A_0=100
\]

\noindent tra un buffer di liquidit\`a \(x_1\) e un'esposizione LDI \(x_2\):

\[
x_1+x_2=100,
\qquad
x_1\geq 0,
\qquad
x_2\geq 50.
\]

I rendimenti assegnati sono rispettivamente

\[
0.010
\qquad\text{e}\qquad
0.050.
\]

Dopo la decisione iniziale, si realizza uno tra tre scenari

\[
\mathcal{S}
=
\{N,S,E\},
\]

\noindent con probabilit\`a

\[
p_N=0.50,
\qquad
p_S=0.30,
\qquad
p_E=0.20.
\]

Il margin call nello scenario \(s\) dipende dall'esposizione LDI:

\[
m_s
=
\alpha_s x_2,
\]

\noindent con

\[
\alpha_N=0.10,
\qquad
\alpha_S=0.25,
\qquad
\alpha_E=0.40.
\]

Dopo l'osservazione dello scenario, il fabbisogno pu\`o essere fronteggiato
mediante liquidit\`a esterna \(e_s\) e liquidazione della posizione LDI
\(f_s\). La liquidazione non produce necessariamente liquidit\`a per l'intero
ammontare ceduto: nello scenario \(s\), il coefficiente di realizzo \`e
\(\beta_s\).

Il bilancio di liquidit\`a \`e

\[
x_1
+
e_s
+
\beta_s f_s
=
\alpha_s x_2
+
g_s,
\]

\noindent dove \(g_s\geq 0\) rappresenta l'eventuale liquidit\`a residua.

Le decisioni di ricorso sono inoltre soggette ai limiti

\[
0\leq e_s\leq 4,
\]

\[
0\leq f_s\leq 25,
\]

\[
f_s\leq x_2.
\]

Il costo di ricorso \`e

\[
\kappa_s e_s
+
\lambda_s f_s.
\]

Rispetto al Caso Aula, emerge quindi una differenza essenziale. Nel problema
SVB, il funding di emergenza non presenta un limite superiore e garantisce
ricorso fattibile per ogni allocazione iniziale ammissibile. Nel Caso
TakeHome, invece, sia la liquidit\`a esterna sia il deleveraging sono
limitati.

La fattibilit\`a della decisione iniziale deve pertanto essere verificata
scenario per scenario.

Questo punto diventa particolarmente importante nel confronto con la politica
Expected Value. Il problema costruito con i parametri medi pu\`o produrre una
decisione perfettamente ammissibile nel modello medio ma non rivalutabile in
uno degli scenari originari.

Con i parametri assegnati nella Scheda Caso, la soluzione del problema medio
\`e

\[
x^{EV}
=
(0,100)^{\prime}.
\]

Nello scenario estremo \(E\), il margin call richiesto \`e

\[
m_E
=
0.40\cdot 100
=
40.
\]

La liquidit\`a massima ottenibile mediante il ricorso disponibile \`e invece

\[
4
+
0.75\cdot 25
=
22.75.
\]

Poich\'e

\[
22.75<40,
\]

\noindent la decisione \(x^{EV}\) non ammette un ricorso fattibile nello
scenario estremo.

Il risultato non deve essere corretto introducendo autonomamente nuove
variabili o nuove fonti di liquidit\`a. La Scheda Caso definisce infatti la
specifica del problema e i limiti di ricorso ne fanno parte.

Il controllo di fattibilit\`a precede quindi il calcolo di qualunque indicatore
di confronto. Se la decisione Expected Value non pu\`o essere rivalutata negli
scenari originari, il benchmark EV non \`e disponibile come valore finito
nella forma utilizzata nel Caso Aula.

Il Caso TakeHome mostra cos\`i che

\[
\text{soluzione del problema medio}
\not\Rightarrow
\text{fattibilit\`a negli scenari originari}.
\]

Questa criticit\`a rende operativa una distinzione gi\`a introdotta nei
Capitoli 14 e 15: la presenza o meno di ricorso relativamente completo.

Il problema TakeHome non costituisce quindi una semplice replica del Caso
Aula con parametri differenti. La limitazione delle azioni di ricorso modifica
la struttura della verifica e rende la fattibilit\`a stessa parte essenziale
dell'analisi.

\section{Sintesi del capitolo}

Il capitolo ha tradotto la programmazione stocastica a due stadi in una
procedura computazionale basata su \texttt{scipy.optimize.linprog}.

La sequenza fondamentale \`e

\[
\text{decisione ex ante}
\longrightarrow
\text{scenario}
\longrightarrow
\text{ricorso}
\longrightarrow
\text{valore della politica}.
\]

La soluzione stocastica utilizza congiuntamente tutti gli scenari; il problema
Expected Value sostituisce l'incertezza con parametri medi; la rivalutazione
EV riporta la decisione ottenuta negli scenari originari; il benchmark
Wait-and-See misura invece il valore di una situazione nella quale lo scenario
sia noto prima della decisione.

Quando i benchmark sono ben definiti e fattibili, il confronto tra

\[
z^{EV},
\qquad
z^{SP},
\qquad
z^{WS}
\]

\noindent permette di costruire \(VSS\) ed \(EVPI\) e di valutare il costo
economico di differenti strutture informative.

I due casi hanno inoltre mostrato che il controllo non costituisce una fase
accessoria della soluzione. Deve essere verificato che la decisione iniziale
sia trattata coerentemente, che il ricorso rispetti i limiti assegnati e che
ogni benchmark rimanga fattibile nel problema nel quale viene valutato.

Il passaggio conclusivo pu\`o quindi essere rappresentato come

\begin{gather*}
\text{modello}\\
\downarrow\\
\text{problema computazionale}\\
\downarrow\\
\text{soluzione}\\
\downarrow\\
\text{controllo}\\
\downarrow\\
\text{interpretazione}.
  \end{gather*}

Python rende operativo il modello; il solver produce una soluzione numerica;
la verifica stabilisce se quella soluzione possa essere utilizzata come base
per una decisione finanziaria.
